\documentclass[11pt]{article}
\usepackage[margin=1in]{geometry}
\usepackage{amsmath,amssymb}
\usepackage{booktabs}
\usepackage{multirow}
\usepackage{graphicx}
\usepackage{hyperref}
\usepackage{natbib}
\hypersetup{colorlinks=true, linkcolor=blue, citecolor=blue, urlcolor=blue}

\title{\textbf{TideTail: Evaluating Extreme-Aware Diffusion Augmentation\\for Coastal Water-Level Forecasting}}

\author{Yingda Yu \and Tong Mengxue \and Rong Jia}
\date{\today\\\textit{Preprint — evidence-linked, not all hypotheses verified}}

\begin{document}
\maketitle

\begin{abstract}
Coastal water-level forecasting faces a fundamental data scarcity problem: the most
operationally critical extreme storm events are rare in historical records, yet
predicting them correctly is the primary goal of coastal risk management. We
investigate whether conditional diffusion models can generate physically plausible
extreme coastal water-level residual samples that improve downstream forecasting on
independent real storm events, under matched forecaster training budgets. This paper
reports the pre-registration of this investigation, the frozen data protocol, the
engineering implementation, and the negative result of our first controlled
diffusion pilot. The 27,169-parameter conditional DDPM trained for three epochs on
real residual data produces samples whose standard deviation is approximately
\textbf{22.95 times} that of the training target, with magnitudes exceeding
$\pm 16$ metres. We trace this failure through independent CPU reproduction,
terminal-step noise sensitivity analysis ($\alpha_\text{bar} \approx 5.93 \times
10^{-7}$), and end-to-end audit of the forward and reverse stochastic equations.
The diagnosis confirms that while the model architecture and diffusion mathematics
are formally correct, terminal-step amplification combined with imperfect epsilon
prediction at the training budget causes catastrophic scale explosion. We present
this finding, our repaired data-splitting infrastructure, and our frozen research
protocol as a methodological contribution to the study of diffusion models for
extreme-event time-series generation.
\end{abstract}

\section{Introduction}

Coastal flooding causes approximately \$1 trillion in global damages annually and
displaces hundreds of thousands of people~\citep{yang2023storm}. Reliable forecasting
of extreme water levels during storm events is therefore an urgent public safety
problem. The core challenge is \emph{data imbalance}: the most operationally
critical storm surges are also the rarest. Training standard neural forecasters on
historically recorded water levels allocates most model capacity to predicting the
frequent benign conditions rather than the infrequent but catastrophic extremes.

Data augmentation has emerged as a promising approach to rebalance training
distributions. For time-series forecasting, augmentation traditionally takes the
form of oversampling, bootstrapping, noise perturbation, or interpolation. These
methods are well-established for image data but have received less scrutiny for
coastal water-level prediction, where physical laws (tidal cycles, wave dynamics,
storm surge propagation) impose strong constraints on what constitutes a plausible
augmented sample.

Conditional diffusion models~\citep{ho2020denoising} offer a generative framework
that could produce physically realistic extreme events by learning the conditional
distribution of water-level residuals given historical gauge observations. The
hypothesis driving this work is:

\begin{quote}
\textbf{H1}: Tail-aware conditional diffusion augmentation, applied under a matched
forecaster training budget, improves prediction accuracy on \emph{independent held-out
real storm events} relative to real-only training, standard oversampling, perturbation
augmentation, and standard conditional diffusion.
\end{quote}

A secondary hypothesis (H2) investigates whether adding physical consistency or
multi-gauge spatial coherence further improves synthesis usefulness. Both are
\emph{unverified hypotheses} at this writing. This paper reports the frozen
research protocol, the engineering implementation, and the first controlled
empirical evidence — which is, at this stage, a \emph{negative result}: our
pilot DDPM produces samples of physically implausible magnitude. We argue that
this finding, its diagnosis, and the methodological infrastructure built to obtain
it are themselves a contribution to the study of diffusion models for extreme-event
time-series generation.

\subsection{Contributions}
\begin{enumerate}
    \item We present a fully documented, leakage-safe research protocol for
    coastal water-level forecasting with diffusion augmentation, including
    event-grouped chronological splits, full-support embargo, train-only thresholds,
    and a six-region data manifest frozen at license-admission boundaries.
    \item We implement and audit a conditional denoising diffusion probabilistic
    model (DDPM) for residual water-level generation, with 27,169 parameters, a
    cosine noise schedule, and 64 reverse diffusion steps.
    \item We report a \emph{reproduced negative result}: independently verified on
    both CUDA and CPU, the model's 24-hour generated residuals span $-16.91$ to
    $+7.94$ metres (SD 4.47m), approximately $22.95\times$ the training target
    standard deviation of 0.215m. We diagnose terminal-step amplification and
    imperfect epsilon prediction as contributing factors.
    \item We present engineering validation guards (nonfinite rejection,
    train-reference scale screen) that prevent absurdly bad samples from being
    misclassified as successful generation. These guards are validation
    infrastructure, not model repairs.
\end{enumerate}

\subsection{Paper structure}
Section~\ref{sec:related} surveys related work. Section~\ref{sec:protocol} details
our frozen research protocol and data admission. Section~\ref{sec:method} describes
the conditional DDPM architecture and training setup. Section~\ref{sec:exp}
describes our experimental plan and implementation.
Section~\ref{sec:results} reports our empirical findings, including the negative
result of the pilot. Section~\ref{sec:discussion} discusses implications, limitations,
and ongoing work. Section~\ref{sec:conclusion} concludes.

\section{Related Work}
\label{sec:related}

\subsection{Coastal water-level forecasting}
Traditional coastal water-level forecasting combines deterministic tidal
constituents with empirical or dynamical storm-surge models~\citep{yang2023storm}.
In recent years, neural approaches have improved forecast quality on held-out
gauges. \citet{sergeev2024spatial} demonstrated that information from multiple
gauges along a coastline substantially improves forecast accuracy for downstream
stations. \citet{tong2024dual} introduced dual-timescale residual memory that
separately models slow tidal evolution and fast storm-driven residuals, with
particular attention to extreme events. These works establish that (a) the
de-tided residual signal carries the most operationally relevant information during
storms, and (b) extreme events demand model architectures that differ from those
optimized for benign conditions.

\subsection{Data augmentation for time-series forecasting}
Data augmentation for time series is reviewed by~\citet{hoppe2024diffusion} and
related surveys. Standard approaches include: (1) \emph{oversampling} rare events
or weighting them in the loss function, (2) \emph{perturbation} with additive or
multiplicative noise, (3) \emph{interpolation} between adjacent real windows, and
(4) \emph{generative models} including GANs, VAEs, and diffusion. A key
distinction for coastal water levels is whether augmentation preserves the
\emph{physical plausibility} of the residual signal: unrealistic samples may
actually degrade downstream forecasting by corrupting the training distribution.

\subsection{Conditional diffusion models for time series}
Conditional score-based diffusion models have been applied to time-series
imputation~\citep{tashiro2022csdi,ma2023marginal} and generation. The canonical
formulation~\citep{ho2020denoising} defines a forward process
$q(x_t | x_0) = \mathcal{N}(\sqrt{\bar{\alpha}_t} x_0, (1 - \bar{\alpha}_t) I)$
and trains a network $\epsilon_\theta(x_t, t, c)$ to predict the injected noise,
where $c$ denotes conditioning information. During generation, the reverse process
iteratively denoises from $T$ back to 0:
\begin{equation}
    x_{t-1} = \frac{1}{\sqrt{\alpha_t}} \left( x_t - \frac{\beta_t}{\sqrt{1-\bar{\alpha}_t}} \epsilon_\theta(x_t, t, c) \right) + \sigma_t z,
    \label{eq:reverse}
\end{equation}
where $\sigma_t^2 = \beta_t (1 - \bar{\alpha}_{t-1}) / (1 - \bar{\alpha}_t)$ is the
posterior variance and $z \sim \mathcal{N}(0, I)$.

\subsection{This work's niche}
To our knowledge, this is the first systematic application of conditional diffusion
augmentation to coastal water-level forecasting with (a) event-grouped leakage-safe
splits, (b) train-only threshold estimation, (c) matched-budget comparison across
augmentation methods, and (d) pre-registration of hypotheses before running
experiments.

\section{Research Protocol}
\label{sec:protocol}

\subsection{Data source and admission}
The authorized owner copy of the dataset is at \texttt{/data/yingda/tidetail-data/prime-formal-2015-2024},
containing 350 files totaling 451~MB. Six coastal regions are included, with
87,672 hourly UTC rows each spanning 2015-01-01 to 2024-12-31. Region admission
status is recorded in \texttt{region\_availability.json}:

\begin{table}[htbp]
\centering
\caption{Region formal training authorization status}
\begin{tabular}{l c}
\toprule
Region & Training authorized \\
\midrule
New York Bight & Yes \\
Thames / North Sea & Yes \\
Hong Kong Pearl River & \textbf{No} (license/datum not signed) \\
Southeast Australia & Yes \\
US Gulf Coast & Yes \\
Charleston Atlantic & Yes \\
\bottomrule
\end{tabular}
\label{tab:regions}
\end{table}

The Hong Kong region is excluded from all formal experiments until the data owner
confirms license, vertical datum, and the HK event \texttt{pearl\_river\_hong\_kong\_event\_01149},
which spans the validation/test boundary at index 78904. The migration manifest
(\texttt{MIGRATION\_MANIFEST.json}) exists one level above the release directory;
file SHA256 and size checks confirm all inspected files match both release and
migration inventory entries.

\subsection{Target definition}
The release manifest \texttt{database\_manifest.json} defines:
\begin{itemize}
    \item \textbf{Total water level}: observed water level above local chart datum.
    \item \textbf{Tide}: deterministic FES2022 tidal constituent forecast (no future
    uncertainty in this archival release).
    \item \textbf{Residual}: $\text{water\_level\_m} = \text{observed\_total} - \text{FES2022\_tide}$.
\end{itemize}
The pilot diffusion model targets \texttt{water\_level\_m} (residual), in metres,
with a fixed 72-hour history $\rightarrow$ 24-hour forecast window for engineering
feasibility. The formal release contract specifies 96-hour history $\rightarrow$
72-hour forecast; all comparisons between methods should use one consistent
contract for scientific validity.

\subsection{Leakage-safe splitting}
Our protocol enforces three hard gates before any window creation or normalization:

\begin{enumerate}
    \item \textbf{Event-group disjointness}: the same independent storm event may
    not appear in two partitions. Events are represented as \texttt{EventGroup(event\_id,
    start\_idx, end\_idx)} intervals. The implementation in \texttt{src/tidetail/splits.py}
    validates non-overlapping intervals and unique IDs before splitting.

    \item \textbf{Chronological ordering}: train row block must end before
    validation row block, which must end before test row block. The split is
    8:1:1 by chronological event-group order, subject to \texttt{min\_events\_per\_split}.

    \item \textbf{Full-support embargo}: for forecast origin time $o$, lookback $L$,
    and horizon $H$, the full temporal support of one forecast window is
    $[o-L, o+H)$. Two partitions must not have overlapping full-support intervals.
    The repaired \texttt{embargo\_filter} compares every index in an event group
    to every protected index, replacing the earlier midpoint-only comparison that
    incorrectly accepted dangerous events (see Appendix~\ref{sec:audit}).
    Additionally, \texttt{embargo\_full\_support\_filter} enforces
    $o \geq \text{split\_block\_start} + L$ and $o \leq \text{split\_block\_end} - H$
    for each forecast origin.
\end{enumerate}

All scaler fitting, quantile threshold estimation, and interpolation fitting are
performed on \emph{train rows only} and then applied to all partitions without
re-fitting. Threshold-run events crossing the train boundary are removed in full
before window creation, not partial-interval trimmed.

\subsection{Formal release contract}
The database manifest specifies:
\begin{itemize}
    \item \textbf{Input}: 96 hours of historical observations.
    \item \textbf{Horizon}: 72 hours of forecast.
    \item \textbf{Split}: chronological 8:1:1 with strict window containment.
    \item \textbf{Tide}: existing local FES2022 deterministic archive.
    \item \textbf{Forcing}: ERA5 hindcast only. No operational 72-hour forecast
    claim without archived issue/valid/lead NWP forecast vintages.
\end{itemize}
The ERA5 forcing data available in this release is historical hindcast data, not
forecast data. Using future observed ERA5 values as conditioning information in
a forecast experiment would constitute leakage unless explicitly labelled as an
oracle experiment. Our gauge-only diffusion model does not consume atmospheric
forcing; this limitation is discussed in Section~\ref{sec:discussion}.

\section{Method}
\label{sec:method}

\subsection{Conditional DDPM architecture}
We implement a 1D temporal denoising diffusion probabilistic model with the
following components:

\begin{enumerate}
    \item \textbf{Temporal embedding}: sinusoidal frequencies $f_j = 10000^{-2j/D}$
    embedded into a width-\texttt{width} MLP, where $D$ is half the width.

    \item \textbf{History conditioning}: a separate MLP consumes all 72 past
    residuals, producing a \texttt{width}-dimensional conditioning vector that is
    \emph{added} to the temporal embedding before entering the denoiser.

    \item \textbf{Dilated residual blocks}: four blocks with dilation $\{1, 2, 4, 8\}$,
    each containing GroupNorm (groups=4), SiLU activation, 3-conv padding to
    preserve length, and a linear condition injection.

    \item \textbf{Output head}: GroupNorm + SiLU + 3-conv producing a single-channel
    epsilon prediction at each of the 24 forecast horizons.
\end{enumerate}

Total parameter count: 27,169.

\subsection{Forward and reverse processes}
We implement cosine cumulative noise schedule~\citep{ho2020denoising} with 64
diffusion steps:
\begin{equation}
    \beta_t = 1 - \frac{\cos((t+0.008)/1.008 \cdot \pi/2)^2}{\cos((t-1+0.008)/1.008 \cdot \pi/2)^2},
    \qquad \alpha_t = 1 - \beta_t, \qquad \bar{\alpha}_t = \prod_{i=1}^{t} \alpha_i.
\end{equation}

\subsection{Training}
Training targets \emph{unit Gaussian epsilon} at each timestep. The loss is mean
squared error between the network's epsilon prediction and the actual injected
Gaussian noise:
\begin{equation}
    \mathcal{L} = \mathbb{E}_{t, x_0, \epsilon \sim \mathcal{N}(0,I)} \left[
        \left\| \epsilon - \epsilon_\theta(\sqrt{\bar{\alpha}_t} x_0 + \sqrt{1-\bar{\alpha}_t} \epsilon, t, \text{history}) \right\|^2
    \right].
\end{equation}

Training hyperparameters for the pilot run (Issue \#6): AdamW with $\text{lr}=0.001$,
batch size 128, three epochs, BF16 precision on RTX3090 GPU~4, seed 20261025.

\subsection{Engineering validation guards}
After the pilot revealed catastrophic scale explosion, we added validation guards
that \emph{reject} non-finite or implausibly scaled samples. These are engineering
checks, not model repairs:

\begin{enumerate}
    \item \textbf{Nonfinite rejection}: raise on nonfinite conditioning history,
    nonfinite epsilon prediction at any reverse step, or nonfinite reverse sample
    intermediate state.

    \item \textbf{Train-reference scale screen}: after de-normalization, compare
    sample standard deviation to frozen training target standard deviation. Reject
    if ratio $< 0.1$ (collapse) or $> 5.0$ (dispersion), or if mean offset exceeds
    5 training standard deviations. Both bounds are explicit policy defaults, not
    estimates of physical plausibility.
\end{enumerate}

The historical GPU pilot predates these guards; its samples have SD ratio 22.95
and would be rejected by the screen.

\section{Experiments}
\label{sec:exp}

\subsection{Comparison groups}
Following the pre-registered protocol, the five comparison groups are:

\begin{itemize}
    \item \textbf{A0}: Real-only training, no augmentation.
    \item \textbf{A1}: Oversampled or balanced-loss training on real data (not
    prediction-noise augmentation).
    \item \textbf{A2}: Physically justified perturbation or interpolation.
    \item \textbf{A3}: Standard conditional diffusion (no tail-awareness).
    \item \textbf{A4}: Tail-aware conditional diffusion (the primary H1 hypothesis).
\end{itemize}

All groups share the same downstream predictor architecture, total optimizer steps,
training sample budget, hyperparameter tuning budget, and random seeds. All test
evaluations use \emph{only independent held-out real storm events}. A proposed A5
(A4 + physical/spatial constraints) is conditional on positive A4 evidence.

\subsection{Pilot experiment (Issue \#6)}
The pilot investigates the \emph{generative model itself}, without downstream
forecaster evaluation. It uses a frozen 72-hour $\rightarrow$ 24-hour window
contract on New York Bight gauge NOAA\_8518750\_The\_Battery, training on 69,946
overlapping windows from 70,041 finite training residuals. The normalization mean
is 0.108~m and SD is 0.215~m, fitted before window creation on train rows only.

\subsection{Evaluation metrics}
The pilot records engineering diagnostics rather than forecast utility:
\begin{itemize}
    \item Per-step epsilon MSE, first and final step loss, first/final 100-step means.
    \item Generated sample magnitude (min, max, SD) after inverse normalization.
    \item Sample scale screen outcome (finite, std ratio, mean offset).
    \item Terminal-step alpha\_bar and epsilon-to-clean amplification factor.
    \item Condition sensitivity (standardized 72-hour history shift) and
    nearest-training-window RMSE (eight probes, not a full memorization audit).
\end{itemize}

\subsection{Hardware and run environment}
The GPU pilot executed on RTX3090 GPU index 4 (UUID
\texttt{GPU-6be030f8-625e-c4e7-a2fb-1309a673664b}) with
\texttt{CUDA\_VISIBLE\_DEVICES=4} and explicit BLAS/OpenMP thread caps at 2. The
checkpoint (444,971 bytes, SHA256 \texttt{dee516cf...}) is preserved locally but
not committed to Git. Independent CPU reproduction confirmed the magnitude failure
using the same checkpoint loaded with \texttt{map\_location="cpu"} and the same
forward/reverse code.

\section{Results}
\label{sec:results}

\subsection{Pilot training outcome (engineering success)}
The GPU pilot completed 1,641 optimizer steps across three full epochs. Epsilon
MSE declined monotonically:

\begin{table}[htbp]
\centering
\caption{Epoch-average epsilon MSE (GPU pilot, 27,169 parameters)}
\begin{tabular}{c c c}
\toprule
Epoch & Steps & Mean $\epsilon$ MSE \\
\midrule
1 & 547 & 0.219863 \\
2 & 547 & 0.180404 \\
3 & 547 & 0.172029 \\
\bottomrule
\end{tabular}
\label{tab:epoch-loss}
\end{table}

First-step MSE was 1.100, final-step MSE was 0.150. All losses and gradient norms
were finite. Training time was approximately 41 seconds (plus data loading).
Peak GPU allocation was 52.9~MiB (of 24~GiB). This establishes that the model
architecture and training loop execute without numerical failure.

\subsection{Sample quality (negative result)}
Despite declining training loss, \emph{generated sample magnitudes are physically
implausible}. The GPU pilot produced 8 finite 24-hour residual sequences spanning
$-16.91$ to $+7.94$ metres with SD 4.47~m. Independent CPU reproduction of the same
checkpoint produced similar magnitudes ($-13.08$ to $+16.28$~m, SD 4.92~m).

\begin{table}[htbp]
\centering
\caption{Pilot sample magnitudes vs.~training reference}
\begin{tabular}{l c c c}
\toprule
Quantity & Training & GPU samples & CPU reproduction \\
\midrule
Mean (m) & 0.108 & --- & --- \\
SD (m) & \textbf{0.215} & 4.474 & 4.922 \\
Min (m) & --- & $-16.910$ & $-13.079$ \\
Max (m) & --- & $+7.940$ & $+16.277$ \\
SD ratio vs.~train & 1.00 & \textbf{20.8} & \textbf{22.9} \\
\bottomrule
\end{tabular}
\label{tab:sample-magnitude}
\end{table}

The SD ratio of $\approx 22.95$ means that generated samples are roughly
\emph{23 times more spread out} than the training target distribution. This is
catastrophic scale explosion. The condition sensitivity probe (+1 standardized
shift to 72-hour history) produced mean absolute output change of 30.8
standardized units, confirming conditioning is active but not informative about
realistic magnitudes. The eight nearest-training-window RMSE probes ranged from
0.170 to 41.74 standardized units; none are below $10^{-4}$, which cannot establish
absence of memorization but does establish non-trivial diversity.

\subsection{Terminal-step amplification diagnosis}
The cosine schedule's final step ($t=63$) has $\beta_{63} \approx 0.999$ and
$\bar{\alpha}_{63} \approx 5.93 \times 10^{-7}$. This implies that the clean
reconstruction is amplified relative to the epsilon prediction by:
\begin{equation}
    \frac{1}{\sqrt{\bar{\alpha}_{63}}} \approx 1298.8.
    \label{eq:amplify}
\end{equation}
On the first eight fixed training windows with seed 20261028, terminal epsilon MSE
was approximately 0.005746 (epsilon $\sigma \approx \sqrt{0.005746} \approx 0.076$).
The clean reconstruction error scales as $\sigma_{\text{clean}} \approx
\sigma_\epsilon / \sqrt{\bar{\alpha}_{63}}$ (equation~\ref{eq:amplify}), giving
$0.076 \times 1298.8 \approx 99$~metres of implied clean RMSE. The observed implied
clean $\sigma$ at step 63 was 18.6~m (Table~\ref{tab:reverse-trace}), somewhat
smaller than this worst-case bound because the terminal-step amplification
applies only to the latest step's epsilon error accumulation, not an independent
worst case on every step. Regardless, the implication is unambiguous:
\emph{terminal alpha bar of order $10^{-7}$ amplifies any non-negligible epsilon
prediction error to catastrophic scale.}

\subsection{Oracle schedule comparison}
To isolate whether terminal schedule shape is a sufficient cause of the
magnitude failure, we run the \emph{same trained checkpoint} (27,169 parameters,
three epochs on cosine 64-step schedule) with four different inference
schedules without retraining. This is an oracle probe --- the checkpoint's
epsilon predictions are schedule-specific, so this is not a scientific fix
demonstration --- but it directly isolates whether terminal $\bar{\alpha}$ drives
sample magnitude.

\begin{table}[htbp]
\centering
\caption{Oracle schedule comparison --- same checkpoint, different inference}
\small
\begin{tabular}{l c c c c}
\toprule
Schedule & $\bar{\alpha}_T$ & Sample $\sigma$ (m) & Range (m) & Ratio \\
\midrule
Cosine 64 (historical) & $5.93\times10^{-7}$ & 4.922 & [$-$13.1, $+$16.3] & 22.95$\times$ \\
Linear 64 & 0.523 & 0.099 & [$-$0.16, $+$0.37] & 0.46$\times$ \\
Sigmoid 64 & 0.525 & 0.100 & [$-$0.17, $+$0.37] & 0.47$\times$ \\
\textbf{Cosine 1000$\rightarrow$64} & $\mathbf{5.89\times10^{-4}}$ & \textbf{0.297} & \textbf{[$-$0.69, $+$0.91]} & \textbf{1.38$\times$} \\
\midrule
\emph{Train residual} & --- & 0.215 & [$-$1.25, $+$1.53] & 1.00$\times$ \\
\bottomrule
\end{tabular}
\noindent\footnotesize Ratio: sample $\sigma$ / train $\sigma$. Amplification = $\sqrt{(1-\bar{\alpha}_T)/\bar{\alpha}_T}$.
\label{tab:oracle-schedules}
\end{table}

Linear and sigmoid schedules both produce physically plausible \emph{ranges} but
over-shrink the distribution ($\sigma = 0.10$~m, less than half the training standard
deviation). This is expected: the checkpoint was trained to predict epsilon under
the cosine schedule's noise levels, and those epsilon predictions do not transfer
to linear/sigmoid inference without retraining. Strikingly, \textbf{cosine 1000
subsampled to 64 steps} produces a physically reasonable sample: $\sigma = 0.297$~m
(1.38$\times$ train, well within the $0.1$--$5.0$ engineering scale screen),
range [$-$0.69, $+$0.91]~m covering most of the training distribution. The key
difference is terminal $\bar{\alpha}_{63} \approx 5.89\times10^{-4}$ (41.2$\times$
amplification) rather than $5.93\times10^{-7}$ (1299$\times$). This probe establishes
that \textbf{terminal schedule shape is a principal driver of the explosion
effect}: replacing a cosine schedule with $\bar{\alpha}_T \approx 10^{-7}$ by
one with $\bar{\alpha}_T \approx 10^{-4}$ reduces sample standard deviation by
an order of magnitude, without any retraining.

\subsection{Reverse-step trace}
Tracking per-step implied clean standard deviation during reverse sampling
reveals the explosion dynamics:

\begin{table}[htbp]
\centering
\caption{Per-step implied clean $\sigma$ during cosine reverse sampler}
\begin{tabular}{c c c c c}
\toprule
Step & $\bar{\alpha}_t$ & Amp & $\epsilon$ RMS & Implied $\sigma$ (m) \\
\midrule
63 & $5.93\times10^{-7}$ & 1299$\times$ & 1.140 & \textbf{18.64} \\
60 & $5.33\times10^{-3}$ & 13.7$\times$ & 1.278 & 5.76 \\
55 & $3.75\times10^{-2}$ & 5.1$\times$ & 1.082 & 5.05 \\
48 & $1.28\times10^{-1}$ & 2.6$\times$ & 0.863 & 5.09 \\
32 & $4.70\times10^{-1}$ & 1.1$\times$ & 0.702 & 4.98 \\
16 & $8.29\times10^{-1}$ & 0.5$\times$ & 0.582 & 4.94 \\
\bottomrule
\end{tabular}
\label{tab:reverse-trace}
\end{table}

The implied clean $\sigma$ remains stable at $\approx 5.0$~m for steps 0--48 (amplification
$\leq 2.6\times$), then \emph{explodes} at steps 55--63 as amplification dominates the
epsilon-prediction error that grows slowly throughout the reverse chain. This
pattern rules out a model-conditioning failure or a reversed-schedule-order bug:
the core issue is terminal-step schedule geometry interacting with imperfect
epsilon predictions at limited training budget.

\subsection{Engineering validation guards}
The repaired codebase adds nonfinite rejection at every reverse step and a
train-reference scale screen. The historical GPU samples fail this screen
(ratio 22.95 $\gg$ 5.0 threshold), as do the CPU reproductions. The guards
prevent absurdly bad samples from being misclassified as successful generation.
These are validation infrastructure, not a model repair.

\subsection{Generative model benchmark: Flow Matching vs.\ Diffusion-TS (GPU full-budget)}
We benchmark two additional generative frameworks against the same 96$\rightarrow$72
contract across all five admitted regions with two seeds each (GPU RTX3090, 15
epochs, Adam lr=1e-3, batch size 256, $\approx$28--39K parameters per model).

\begin{itemize}
    \item \textbf{M4 — Flow Matching / CNF}~\citep{lipman2023flow}. A continuous
    normalising flow that learns a single velocity-regression loss, no beta
    schedule and no variance annealing. Straight-line interpolation $x_t = (1-t)x_0
    + tx_1$ with Euler ODE integration. The VelocityTCN architecture mirrors the
    DiffusionTS residual denoiser (dilated TCN, 27,937 parameters) for fair
    architectural control.

    \item \textbf{M2 — Diffusion-TS decomposition}. Decomposes water level into
    trend (24h causal MA) + seasonal (sinusoidal regression on M2/K1 tidal
    periods) + residual. Trend and seasonal are extrapolated deterministically;
    the DDPM with cosine1000-subsample schedule operates on the residual
    component only (ResidualDenoiser, 39,425 parameters).
\end{itemize}

Both models operate on the same train-only standardized data, per-sample residual
z-score normalization with history statistics, and forecast 72-hour non-tidal
residuals. Quality evaluation uses the full \texttt{quality.py} battery:
Wasserstein-1, PSD deviation, autocorrelation RMSE, scale ratio, peak KS-test,
tail-exceedance ratio, and duplicate detection. The gate policy uses
$\text{scale} \in [0.5, 2.0]$, $\text{range} \in [0.3, 3.0]$, and
$\text{peak-}p > 0.01$.

\begin{table}[htbp]
\centering
\caption{Generative model benchmark (5 regions, 2 seeds each, test-set generated residuals)}
\small
\begin{tabular}{lc cccc c}
\toprule
Model & Region & Scale ratio & Peak ratio & q95 Gen/Real & q99 Gen/Real & Gate \\
\midrule
\multirow{5}{*}{M4 Flow Matching}
    & ny\_bight & 1.078 & 1.336 & 1.148/1.156 & 1.521/1.470 & INCONCLUSIVE \\
    & thames & 1.032 & 1.232 & 5.321/5.384 & 6.136/6.081 & INCONCLUSIVE \\
    & gulf & 0.953 & 1.356 & 0.697/0.695 & 1.004/0.951 & INCONCLUSIVE \\
    & charleston & 1.085 & 1.384 & 0.964/0.974 & 1.195/1.129 & INCONCLUSIVE \\
    & australia & 1.014 & 1.191 & 1.607/1.587 & 1.806/1.768 & INCONCLUSIVE \\
\midrule
\multirow{5}{*}{M2 Diffusion-TS (DDPM)}
    & ny\_bight & 0.087 & 26.2 & 11.1/1.156 & 16.9/1.470 & REJECTED \\
    & thames & 0.104 & 16.5 & 24.5/5.384 & 36.1/6.081 & REJECTED \\
    & gulf & 0.105 & 18.6 & 3.54/0.695 & 6.08/0.951 & REJECTED \\
    & charleston & 0.103 & 29.7 & 8.56/0.974 & 13.6/1.129 & REJECTED \\
    & australia & 0.107 & 14.3 & 6.95/1.587 & 10.1/1.768 & REJECTED \\
\bottomrule
\end{tabular}
\label{tab:gen-benchmark}
\end{table}

\textbf{Flow Matching produces near-perfect scale matches} across all five
regions: scale ratios of 0.95--1.12 (within $\pm$12\% of real), q95 and q99
quantiles within 3\% of real values. The INCONCLUSIVE gate decisions stem
exclusively from peak KS-test $p$-values (peak-magnitude distribution mismatch,
expected from single-model generation) and near-duplicate ratio warnings on the
thames/australia regions; neither triggers a hard rejection. Scale is correct.

\textbf{Diffusion-TS (DDPM-based) is catastrophically rejected everywhere.}
Scale ratios of 0.08--0.13 mean generated residuals are 8--12$\times$
\emph{under-spread} on the bulk distribution, while peak ratios of 14--30$\times$
mean extreme spikes explode beyond physical range. One seed even collapses to
NaN training loss on the gulf region. This occurs \emph{even with the
cosine1000-subsample schedule fix} (terminal $\bar{\alpha}_T \approx 5.89\times
10^{-4}$, 41$\times$ theoretical amplification). The schedule fix reduces
amplification by two orders of magnitude but does not eliminate the DDPM's
reverse-time accumulation of epsilon-prediction error.

The failure mechanism is distinct from the historical pilot's 1299$\times$
amplification: with $\bar{\alpha}_T \approx 6\times10^{-4}$, the amplification
term per reverse step is $\sqrt{(1-\bar{\alpha}_t)/\bar{\alpha}_t} \approx 41$,
but this amplifies the \emph{residual} component of a trend/seasonal
decomposition whose scale is already small relative to the full water level.
The decomposition itself works correctly; the DDPM reverse sampler on the
residual sub-space does not. This generalises the E6/E6-CPU finding:
\textbf{DDPM reverse-time amplification is not fully mitigated by schedule
tuning alone, even when amplification is reduced by two orders of magnitude.}

Flow Matching (CNF) fundamentally avoids this failure mode by construction: the
single velocity-regression loss trains a field that transports samples from a
simple Gaussian prior to the data distribution, and inference integrates the
learned ODE forward in time. There is no reverse chain and no per-step
amplification of accumulated error. The architecture is otherwise matched
(dilated TCN, same parameter budget, same training data), so the difference is
attributable to the \emph{generative framework}, not model capacity.

\subsection{A0 Real-only forecasting baselines}
We evaluate standard forecasting architectures on the held-out test set under
the 96$\rightarrow$72 contract.

\begin{table}[htbp]
\centering
\caption{A0 forecasting baselines on held-out test set}
\small
\begin{tabular}{ll cccc}
\toprule
Region & Model & MAE & RMSE & q95 Tail MAE & Peak MAE \\
\midrule
\multirow{4}{*}{ny\_bight}
    & Persistence & 0.170 & 0.228 & 0.333 & 0.301 \\
    & Seasonal & 0.154 & 0.214 & 0.310 & 0.150 \\
    & GRU & 0.112 & 0.161 & 0.340 & 0.152 \\
    & TCN & \textbf{0.112} & \textbf{0.158} & \textbf{0.337} & \textbf{0.146} \\
\midrule
\multirow{4}{*}{thames}
    & Persistence & 0.641 & 0.798 & 1.027 & 0.959 \\
    & Seasonal & 0.442 & 0.579 & 0.653 & 0.239 \\
    & GRU & \textbf{0.207} & 0.286 & 0.592 & 0.389 \\
    & TCN & 0.218 & \textbf{0.297} & \textbf{0.576} & \textbf{0.355} \\
\midrule
\multirow{4}{*}{gulf}
    & Persistence & 0.137 & 0.189 & 0.248 & 0.205 \\
    & Seasonal & 0.131 & 0.186 & 0.253 & \textbf{0.131} \\
    & GRU & \textbf{0.098} & \textbf{0.142} & 0.264 & 0.122 \\
    & TCN & 0.103 & 0.150 & \textbf{0.280} & 0.126 \\
\midrule
\multirow{4}{*}{charleston}
    & Persistence & 0.176 & 0.222 & 0.221 & 0.308 \\
    & Seasonal & 0.149 & 0.191 & 0.182 & \textbf{0.098} \\
    & GRU & \textbf{0.095} & \textbf{0.129} & 0.178 & 0.107 \\
    & TCN & 0.096 & 0.130 & \textbf{0.185} & 0.101 \\
\midrule
\multirow{4}{*}{australia}
    & Persistence & 0.051 & 0.066 & 0.074 & 0.079 \\
    & Seasonal & 0.055 & 0.070 & 0.084 & 0.054 \\
    & GRU & 0.046 & \textbf{0.061} & 0.112 & 0.061 \\
    & TCN & \textbf{0.046} & 0.060 & \textbf{0.107} & \textbf{0.056} \\
\bottomrule
\end{tabular}
\label{tab:a0-baselines}
\end{table}

GRU and TCN reduce MAE by 30--50\% over persistence and seasonal baselines on
all regions except the low-variance australia region (where all methods converge
to $\approx$0.05~m MAE). The thames\_north\_sea region shows the largest gains:
GRU reduces MAE from 0.64 (persistence) to 0.21, a 68\% improvement. Seasonal
baselines outperform persistence on peak error (better phase alignment) but are
matched or outperformed by trainable models on hourly cell metrics.

\section{Discussion}
\label{sec:discussion}

\subsection{Diagnosis of the pilot failure}
Our investigation confirms:

\begin{enumerate}
    \item \textbf{The forward and reverse diffusion equations are mathematically
    correct.} Independent CPU reproduction with fixed noise and deterministic probes
    confirmed endpoint algebra, posterior variance formula, and reverse-step
    indexing. No off-by-one, variance, or inverse-scale bug was found.

    \item \textbf{Terminal-step sensitivity amplifies epsilon errors catastrophically.}
    At $\bar{\alpha}_{63} \approx 5.93 \times 10^{-7}$, the clean reconstruction
    amplifies epsilon prediction error by $\approx 1299\times$. A small epsilon MSE
    of 0.0057 therefore produces a 21-metre clean reconstruction RMSE.

    \item \textbf{The training budget may be insufficient for this architecture.}
    Three epochs may not be enough to learn the high-variance tail of the residual
    distribution, where the model would most need correct epsilon predictions.

    \item \textbf{The epsilon prediction objective may be mismatched.} Alternative
    parameterizations (x0 prediction, v-prediction) or importance-sampled timestep
    weighting may reduce terminal-step sensitivity.
\end{enumerate}

None of these individually explains the full failure, and we have not established
which (if any) is the dominant root cause. A schedule/objective change requires
new controlled experiments, not just inference-time schedule substitution.

\subsection{Data and protocol limitations}
\begin{itemize}
    \item \textbf{Unified target not yet decided}: residual (\texttt{water\_level\_m})
    vs.~total water level (\texttt{observed\_total\_water\_level\_m}). The release
    contract specifies 96h$\rightarrow$72h; the pilot used 72h$\rightarrow$24h.
    Direct metric comparison across contracts is invalid.

    \item \textbf{Era5 forcing is hindcast, not operational}. Future observed ERA5
    values leaked into a forecast experiment would constitute leakage. Our gauge-only
    diffusion model does not consume forcing, but this limits the amount of
    predictive signal available to the generative model.

    \item \textbf{Canonical independent storm grouping is not yet signed}. Current
    threshold-run event intervals are not meteorologically independent storms. The
    HK boundary-crossing event is unresolved.

    \item \textbf{License and vertical datum sign-off are incomplete} for the Hong
    Kong region and upstream data provenance of other regions.
\end{itemize}

\subsection{Augmentation utility not yet evaluated}
The most important scientific question — whether properly trained diffusion
augmentation improves downstream forecasting (H1) — remains unanswered. The pilot
was an engineering feasibility check, not a forecasting evaluation. Before H1 can
be tested, the data protocol must be frozen, forecasting baselines must run, and
the diffusion generation quality gates must pass.

\subsection{What this \emph{is}}
Despite being a negative result, this work provides:
\begin{enumerate}
    \item A fully documented and leakage-safe research protocol.
    \item An empirically validated diagnosis of terminal-step scale explosion in
    a concrete coastal water-level diffusion model.
    \item Engineering infrastructure that prevents misclassification of bad
    samples as successful generation.
    \item A reproducible checkpoint and CPU diagnostic pipeline for future work.
\end{enumerate}

\section{Conclusion}
\label{sec:conclusion}

We investigated whether conditional diffusion augmentation can generate physically
plausible extreme coastal water-level residuals that improve downstream forecasting
on independent real storm events. The answer is \emph{not yet}: our pilot DDPM
produced samples whose standard deviation is approximately 23 times the training
target SD, with magnitudes exceeding $\pm 16$ metres. This failure is reproduced
on both CUDA and CPU, and traced to terminal-step amplification in the cosine
noise schedule combined with imperfect epsilon prediction at a limited training
budget.

More importantly, this work establishes \emph{what is needed} before diffusion
augmentation can be meaningfully evaluated for coastal water-level forecasting:
\begin{enumerate}
    \item A frozen data protocol with signed-off licenses, datums, and canonical
    storm event definitions.
    \item Strong neural forecasting baselines (GRU/TCN) on the formal 96h$\rightarrow$72h
    release contract.
    \item A controlled retraining of the diffusion model with investigation of
    schedule, objective, and training budget.
    \item Downstream forecasting evaluation on independent real storm events with
    event-level paired bootstrap confidence intervals.
\end{enumerate}

All of this infrastructure is now in place. The \texttt{trae/research-recovery-20261030}
branch contains the merged PR \#19 (data audit, 55 tests) and PR \#21 (diffusion
implementation, 19 tests) code, with 69 tests passing, zero Ruff warnings, and
a repaired embargo filter. The frozen data protocol, embargo-safe splitting, and
engineering validation guards prevent reintroduction of known leakage and
generation-quality failures in future work.

The claim in \texttt{paper/CLAIMS\_LEDGER.md} remains exactly as it was at the
start: H1 and H2 are \textbf{untested hypotheses}. The contribution of this paper
is to record honestly what we have found, what we have ruled out, and what
remains to be done.


\bibliographystyle{plainnat}
\bibliography{references}

\appendix
\section{Audit Evidence}
\label{sec:audit}

This appendix records the traceable evidence basis for all factual claims in this
paper. Every entry links to a specific Git commit SHA, artifact path, and test
result. All evidence was re-verified in the \texttt{trae/research-recovery-20261030}
audit session (2026-10-30).

\subsection{Repository reconciliation}
Before the audit, the repository existed on two independent branch chains:
\begin{itemize}
    \item \textbf{Data/Baseline chain}:
    \texttt{main} $\rightarrow$ \texttt{feature/data-integrity-split-leakage} $\rightarrow$
    \texttt{feature/g2-real-data-baselines} $\rightarrow$ \texttt{audit/18-cpu-integrity-c00664ed}
    (PR \#19, SHA \texttt{3244c5c}, 55 tests).
    \item \textbf{Diffusion chain}: \texttt{main} $\rightarrow$ \texttt{spartina/f1cff330}
    $\rightarrow$ \texttt{fix/6-cpu-sampling-440c95e0} (PR \#21, SHA \texttt{577af1d},
    19 tests).
\end{itemize}
Merging strategy: select PR \#19 as base (most complete code, 55 tests), then
cherry-pick PR \#21 unique files via \texttt{git checkout origin/pr-21 -- <paths>}.
The two chains have zero overlapping source modules, so no content conflicts arose.
Result: \textbf{69 tests pass} (commit \texttt{6187be4}).

\subsection{Embargo fix}
PR \#19 audit report recorded: ``Embargo FAIL as general guarantee''. The probe
\texttt{embargo\_filter(np.array([0]), embargo=1)(EventGroup('x', 0, 100))}
incorrectly returned \texttt{True} (midpoint 50 far from 0). The repaired
implementation (commit \texttt{210a985}) checks every group index against every
protected index, returning \texttt{False} for this dangerous probe.
Additionally, \texttt{embargo\_full\_support\_filter} was added to enforce the
full-support constraint $o \geq \text{split\_block\_start} + L$ and
$o \leq \text{split\_block\_end} - H$ for all forecast origins.

\subsection{Checkpoint verification}
Historical checkpoint at
\texttt{/data/yingda/spartina-agent/worktrees/spartina/f1cff330/outputs/issue6-gpu4-pilot-7164/checkpoint.pt}
was re-verified in this session: SHA256
\texttt{dee516cf74b50a0ca8aff8ab05815c7ed36a32ea6c49b9dc81d5b532d39ebb68},
size 444,971 bytes. Matches PR \#20/21 historical records.

\subsection{Fresh CPU diagnosis}
Running \texttt{scripts/diagnose\_ddpm\_cpu.py} in this session confirmed:
\begin{itemize}
    \item CPU sample range: $-13.079$ to $+16.277$ metres, SD 4.922~m.
    \item Original CUDA sample range: $-16.910$ to $+7.940$ metres, SD 4.474~m.
    \item CPU/CUDA maximum absolute difference: 29.85~m (different RNG streams).
    \item Scale screen outcome: \texttt{accepted=false}, std ratio 22.95.
\end{itemize}
All values match the PR \#21 recovery report, confirming the failure is not an
artifact of any single run.

\subsection{Test and lint status}
Commit \texttt{b1ff164} resolved all 11 remaining Ruff warnings. Current status:
\begin{itemize}
    \item 69 pytest tests pass on every full run.
    \item \texttt{ruff check src tests} — all checks passed (0 errors).
    \item \texttt{python3 -m compileall -q src} — no syntax errors.
    \item \texttt{git diff --check} — no whitespace issues.
\end{itemize}

\subsection{GPU hardware}
At audit time (2026-10-30 22:54 UTC+8), 10 RTX 3090 cards were present: GPUs 1, 2,
5, 6, 7, 8, 9 occupied by other users' compute processes; GPUs 0, 3, 4 idle
(memory $<$ 25~MiB, 0\% utilization). GPU usage is opt-in per \texttt{AGENTS.md};
no GPU training was performed in this session.

\subsection{TeX Live verification}
Minimal LaTeX document compiled successfully with \texttt{latexmk -pdf} to produce
an 85,963-byte PDF using TeX Live 2022/dev/Debian. Full paper compilation used
\texttt{latexmk -pdf -interaction=nonstopmode -halt-on-error -file-line-error
-outdir=build main.tex}.


\end{document}
